% =====================================================================
%  Artículo: Selección de un modelo de proceso de software para una
%  aplicación móvil de pedidos (Sabores de la Abuela Rosa)
%  Compilar con pdfLaTeX (Overleaf o TeX Live/MiKTeX). Referencias IEEE.
% =====================================================================
\documentclass[11pt,a4paper]{article}

% ---------- Idioma y codificación ----------
\usepackage[utf8]{inputenc}
\usepackage[T1]{fontenc}
\usepackage[spanish,es-tabla,es-nodecimaldot]{babel}

% ---------- Tipografía y página ----------
\usepackage{mathptmx}            % Times, estilo de artículo científico
\usepackage{microtype}
\usepackage[margin=2.5cm]{geometry}
\usepackage{setspace}
\setstretch{1.1}

% ---------- Tablas ----------
\usepackage{array}
\usepackage{tabularx}
\usepackage{booktabs}
\usepackage{multirow}
\usepackage{pdflscape}
\usepackage{ragged2e}
\newcolumntype{L}{>{\RaggedRight\arraybackslash}X}
\newcolumntype{C}[1]{>{\centering\arraybackslash}p{#1}}
\newcolumntype{P}[1]{>{\RaggedRight\arraybackslash}p{#1}}

% ---------- Otros ----------
\usepackage{xcolor}
\usepackage{enumitem}
\setlist{nosep,leftmargin=*}
\usepackage[font=small,labelfont=bf,labelsep=period]{caption}
\usepackage{titlesec}
\titleformat{\section}{\normalfont\large\bfseries}{\Roman{section}.}{0.6em}{\MakeUppercase}
\titleformat{\subsection}{\normalfont\normalsize\bfseries}{\Alph{subsection}.}{0.5em}{}
\usepackage{fancyhdr}
\pagestyle{fancy}
\fancyhf{}

\fancyhead[R]{\small\thepage}
\renewcommand{\headrulewidth}{0.4pt}
\usepackage{cite}                 % citas numéricas [1], [2] estilo IEEE
\usepackage[hidelinks]{hyperref}
\usepackage{url}
\urlstyle{same}

% ---------- Datos del artículo ----------
\title{\textbf{Selección fundamentada de un modelo de proceso de software
para una aplicación móvil de pedidos: comparación de los modelos
cascada, iterativo y ágil en un emprendimiento gastronómico}}

\author{%
  Nombre Apellido del estudiante\\[2pt]
  \small Carrera / Asignatura: Ingeniería de Requisitos\\
  \small Universidad --- Quevedo, Ecuador
}
\date{12 de noviembre de 2026}

% =====================================================================
\begin{document}
\maketitle
\thispagestyle{fancy}

% ---------------------------------------------------------------------
\begin{abstract}
\noindent
Este trabajo selecciona el modelo de proceso de software más adecuado para
desarrollar una aplicación móvil de pedidos destinada a \emph{Sabores de la
Abuela Rosa}, un emprendimiento ficticio de comida típica en Quevedo
(Ecuador). Se identificaron los actores del proceso, se definieron siete
criterios de comparación (estabilidad de requisitos, criticidad,
disponibilidad del cliente, tamaño, regulación, plazo y presupuesto, y riesgo
técnico) y se valoró cada modelo con una escala de 1 a 5. El modelo ágil
obtuvo 31 de 35 puntos, frente a 27 del iterativo y 16 de la cascada. Se
recomienda Scrum ligero con sprints de dos semanas, fecha y presupuesto fijos
y alcance variable, priorizado como producto mínimo viable (MVP). Para cubrir
la debilidad del ágil en regulación, los requisitos de la Ley Orgánica de
Protección de Datos Personales (LOPDP) y la conciliación de pagos se
incorporan a la Definición de Hecho desde el primer sprint. Finalmente, se
detallan las actividades centrales y transversales del framework de
Ingeniería de Requisitos (IR) que se ejecutan en cada uno de los siete
sprints.

\medskip
\noindent\textbf{Palabras clave:} modelos de proceso de software, desarrollo
ágil, Scrum, ingeniería de requisitos, aplicación móvil, protección de datos
personales.
\end{abstract}

% ---------------------------------------------------------------------
\section{Introducción}

La elección del modelo de proceso condiciona la organización del trabajo,
el cumplimiento de plazos y la calidad del producto; no existe un modelo
universal, sino uno adecuado a las características de cada proyecto
\cite{otavo2024,mex2025}. Este artículo aplica esa idea a un caso concreto:
una aplicación móvil de pedidos para un emprendimiento familiar.

\emph{Sabores de la Abuela Rosa} (organización ficticia) vende bolones,
tigrillo, encebollado y tortillas de yuca en un local pequeño del centro de
Quevedo. Recibe unos 70 pedidos diarios por una aplicación de mensajería;
la dueña, Doña Rosa Mendoza (persona ficticia), los anota en un cuaderno y
con frecuencia se equivoca de dirección o pierde mensajes en las horas pico.

El producto solicitado debe permitir a los clientes ver el menú del día,
hacer pedidos, elegir retiro en el local o entrega a domicilio y pagar por
transferencia (más adelante, quizá, con tarjeta). El objetivo del trabajo es
comparar los modelos cascada, iterativo y ágil en este escenario,
seleccionar uno de forma fundamentada y describir cómo se ejecuta el
framework de IR en cada iteración.

La Tabla~\ref{tab:datos} resume los datos clave del escenario.

\begin{table}[htbp]
\centering
\caption{Datos clave del escenario.}
\label{tab:datos}
\small
\begin{tabularx}{\textwidth}{@{}P{4.2cm}L@{}}
\toprule
\textbf{Dato} & \textbf{Valor} \\
\midrule
Equipo & 2 desarrolladores a medio tiempo \\
Plazo & 14 semanas, fecha fija (fiestas de la ciudad) y presupuesto cerrado \\
Disponibilidad del cliente & Muy alta: la dueña responde a diario y prueba cada versión \\
Volatilidad esperada & Muy alta: menú diario, ideas nuevas de la dueña y de su hijo en cada reunión \\
Criticidad & Media: pagos por transferencia y datos personales; sin riesgo para la vida \\
Marco normativo o contractual & LOPDP de Ecuador \cite{lopdp2021,reglamento2023}; acuerdo informal de precio fijo \\
Restricción técnica & Android de gama baja y datos móviles intermitentes: app liviana y tolerante a cortes \\
\bottomrule
\end{tabularx}
\end{table}

% ---------------------------------------------------------------------
\section{Marco teórico}

\subsection{Modelos de proceso}

El \textbf{modelo en cascada} organiza el desarrollo en fases secuenciales
(análisis, diseño, implementación, verificación y mantenimiento) y exige
requisitos completos y estables al inicio, porque el cambio posterior es
costoso \cite{otavo2024}. El \textbf{modelo iterativo} construye el sistema
en ciclos sucesivos planificados de antemano, cada uno con una versión
parcial que se revisa con el cliente. Los \textbf{métodos ágiles}, como
Scrum, trabajan en iteraciones cortas (\emph{sprints}) sobre un backlog
priorizado que se reordena continuamente, con participación constante del
cliente y entregas frecuentes de software funcionando \cite{schon2017}.

En la práctica, pocos equipos aplican un método en forma pura: en una
encuesta internacional de 556 casos, menos del 15\,\% trabajaba de manera
exclusivamente ágil o tradicional, y la agilidad dependía más de las
prácticas elegidas que del nombre del método \cite{kuhrmann2022}.

\subsection{Ingeniería de requisitos en contextos ágiles}

En la IR ágil, las actividades de elicitación, documentación, validación,
negociación y gestión no son fases separadas: se repiten en cada iteración
y solo se detalla lo necesario para el siguiente ciclo, apoyándose en
historias de usuario, prototipos y un backlog priorizado \cite{schon2017}.
Entre sus desafíos destacan el descuido de los requisitos no funcionales
\cite{schon2017}, la gestión del conocimiento sobre requisitos que cambian
\cite{kasauli2021} y la especificación de requisitos de privacidad exigidos
por la ley \cite{canedo2022}.

% ---------------------------------------------------------------------
\section{Metodología}

El trabajo siguió cuatro pasos: (1) análisis del escenario e identificación
de actores; (2) definición de criterios de comparación; (3) valoración de
cada modelo por criterio en una matriz y elección del modelo; y
(4) descripción del framework de IR por iteración para el modelo elegido.

\subsection{Actores del proceso}

La dueña concentra la decisión y el conocimiento del negocio; por eso actúa
como \emph{Product Owner}, mientras el hijo y los clientes son fuentes de
requisitos que ella prioriza (Tabla~\ref{tab:actores}).

\begin{table}[htbp]
\centering
\caption{Actores del proceso.}
\label{tab:actores}
\footnotesize
\begin{tabularx}{\textwidth}{@{}P{3.1cm}P{2.0cm}LP{1.5cm}L@{}}
\toprule
\textbf{Actor} & \textbf{Tipo} & \textbf{Interés principal} & \textbf{Influencia} & \textbf{Rol en el proceso} \\
\midrule
Doña Rosa Mendoza (dueña) & Cliente / patrocinadora & Pedidos sin errores, cobrar todo, estar lista para las fiestas & Alta & \emph{Product Owner}: prioriza el backlog, valida cada incremento, decide el alcance \\
Hijo de Doña Rosa (redes sociales) & Stakeholder interno & Promociones, fidelización, difusión en redes & Media & Proponente de funciones; sus ideas entran al backlog y se priorizan con la dueña \\
Clientes finales & Usuarios & Pedir rápido desde Android de gama baja, con datos intermitentes & Media & Fuente de requisitos de usabilidad; participan en pruebas piloto \\
Desarrolladores (2, medio tiempo) & Equipo de desarrollo & Alcance realista, requisitos claros por sprint & Alta (técnica) & Analizan, estiman, construyen y prueban; uno asume además el rol de \emph{Scrum Master} \\
Personal de cocina y despacho & Usuario interno & Ver pedidos claros, con dirección y estado de pago & Baja & Usuario del panel de pedidos; valida el flujo operativo \\
Repartidor(es) & Usuario interno & Dirección exacta, referencia y teléfono del cliente & Baja & Usuario de la información de entrega \\
Banco / entidad de la transferencia & Actor externo & Ninguno directo & Baja & Origen de los comprobantes que deben conciliarse \\
Superintendencia de Protección de Datos Personales & Regulador & Cumplimiento de la LOPDP & Alta (normativa) & Fuente de requisitos no funcionales obligatorios: consentimiento, finalidad, seguridad, derechos del titular \cite{lopdp2021,reglamento2023} \\
\bottomrule
\end{tabularx}
\end{table}

\subsection{Criterios de comparación}

Se usaron siete criterios: los cinco solicitados por la actividad más dos
que el escenario vuelve decisivos, el plazo fijo y el riesgo técnico
(Tabla~\ref{tab:criterios}). Cada modelo se valoró según su adecuación al
criterio con la escala: 1 = muy poco adecuado, 2 = poco adecuado,
3 = aceptable, 4 = adecuado y 5 = muy adecuado.

\begin{table}[htbp]
\centering
\caption{Criterios de comparación.}
\label{tab:criterios}
\footnotesize
\begin{tabularx}{\textwidth}{@{}C{0.7cm}P{3.0cm}LL@{}}
\toprule
\textbf{\#} & \textbf{Criterio} & \textbf{Qué evalúa} & \textbf{Situación en el escenario} \\
\midrule
C1 & Estabilidad de requisitos & Cuánto cambiarán los requisitos durante el proyecto & Muy baja: menú diario, puntos vs.\ promociones, ideas nuevas en cada reunión \\
C2 & Criticidad & Daño posible si el software falla & Media: pérdidas económicas y datos personales expuestos; sin riesgo para la vida \\
C3 & Disponibilidad del cliente & Frecuencia con que el cliente puede participar y validar & Muy alta: responde a diario y prueba cada versión \\
C4 & Tamaño del proyecto y del equipo & Volumen de funcionalidades y personas & Pequeño: 2 desarrolladores a medio tiempo, pocas funciones núcleo \\
C5 & Regulación & Exigencias legales o contractuales de documentación y control & Media: LOPDP obligatoria; contrato informal de precio fijo \\
C6 & Plazo y presupuesto & Rigidez de la fecha y del costo & Fijos: 14 semanas y presupuesto pequeño cerrado \\
C7 & Riesgo técnico & Incertidumbre tecnológica & Media: Android de gama baja, conexión intermitente, conciliación de transferencias \\
\bottomrule
\end{tabularx}
\end{table}

% ---------------------------------------------------------------------
\section{Resultados}

\subsection{Matriz comparativa}

El modelo ágil obtuvo 31 de 35 puntos, frente a 27 del iterativo y 16 de la
cascada; la cascada solo supera a los demás en regulación
(Tabla~\ref{tab:matriz}).

\begin{table}[htbp]
\centering
\caption{Matriz comparativa de los modelos cascada, iterativo y ágil (escala 1--5).}
\label{tab:matriz}
\footnotesize
\begin{tabularx}{\textwidth}{@{}P{2.6cm}LLL@{}}
\toprule
\textbf{Criterio} & \textbf{Cascada} & \textbf{Iterativo} & \textbf{Ágil} \\
\midrule
C1 Estabilidad de requisitos (muy baja) &
\textbf{1} --- exige requisitos congelados al inicio; cada cambio de Doña Rosa obligaría a rehacer fases &
\textbf{3} --- absorbe cambios entre iteraciones, pero con un plan de iteraciones definido de antemano &
\textbf{5} --- el backlog se reprioriza cada sprint; los cambios se esperan y se gestionan \\
\addlinespace
C2 Criticidad (media) &
\textbf{3} --- su control formal sobra para un sistema sin riesgo vital &
\textbf{4} --- permite revisar seguridad y pagos en cada iteración &
\textbf{4} --- suficiente si se fijan criterios de aceptación y pruebas para pagos y datos \\
\addlinespace
C3 Disponibilidad del cliente (muy alta) &
\textbf{2} --- desaprovecha al cliente: solo participa al inicio y en la entrega &
\textbf{4} --- el cliente valida al final de cada iteración &
\textbf{5} --- la dueña prueba cada incremento y actúa como \emph{Product Owner} \\
\addlinespace
C4 Tamaño (pequeño) &
\textbf{2} --- la documentación por fase consume tiempo que el equipo no tiene &
\textbf{4} --- estructura ligera y manejable &
\textbf{5} --- ceremonias mínimas, adecuado para equipos pequeños \\
\addlinespace
C5 Regulación (LOPDP; contrato informal) &
\textbf{4} --- su documentación facilita evidenciar cumplimiento &
\textbf{4} --- permite documentar por iteración &
\textbf{3} --- exige incluir explícitamente requisitos de privacidad en la Definición de Hecho \\
\addlinespace
C6 Plazo y presupuesto fijos &
\textbf{2} --- entrega todo al final; un retraso deja a la dueña sin nada para las fiestas &
\textbf{4} --- hay versiones parciales utilizables &
\textbf{5} --- fecha y costo fijos con alcance variable: siempre existe un producto usable priorizado \\
\addlinespace
C7 Riesgo técnico &
\textbf{2} --- el rendimiento real se descubre tarde &
\textbf{4} --- se prueba en dispositivos reales en cada iteración &
\textbf{4} --- pruebas tempranas con clientes reales en sus teléfonos \\
\midrule
\textbf{Total (máx.\ 35)} & \textbf{16} & \textbf{27} & \textbf{31} \\
\bottomrule
\end{tabularx}
\end{table}

\noindent\textbf{Modelo elegido para el escenario: ágil (Scrum ligero).}

\subsection{Recomendación fundamentada}

Se recomienda Scrum ligero con siete sprints de dos semanas (uno de inicio y
seis de desarrollo), un MVP usable desde la semana 8 y la versión para las
fiestas en la semana 14. Las razones son cinco:

\begin{enumerate}
  \item \textbf{La volatilidad es el factor dominante.} Con el menú cambiando
  a diario y nuevas ideas en cada reunión, un modelo que congela requisitos
  fracasaría o se rehacería varias veces. En proyectos de tres meses, los
  equipos Scrum superaron a los de cascada en gestión del cambio,
  planificación y adaptabilidad \cite{mex2025}.
  \item \textbf{El cliente está disponible.} La participación diaria de la
  dueña es la condición que el ágil necesita y que la cascada desperdicia;
  la participación continua del cliente es un factor crítico de éxito de la
  IR ágil \cite{schon2017}.
  \item \textbf{La fecha es inamovible y el presupuesto, cerrado.} Con
  alcance variable, lo que no alcance en 14 semanas queda fuera, pero lo
  esencial siempre funciona.
  \item \textbf{El equipo es pequeño.} Dos personas a medio tiempo no pueden
  sostener la documentación pesada de la cascada.
  \item \textbf{El riesgo técnico se reduce pronto.} Probar cada incremento
  en teléfonos Android de gama baja revela problemas de peso y conectividad
  desde el primer mes.
\end{enumerate}

El modelo iterativo también sería viable (27 puntos), pero planifica las
iteraciones de antemano y aquí el contenido cambia demasiado para ese plan.
Además, el modelo ágil se adapta en lugar de aplicarse en forma pura, en
línea con la evidencia de que la mayoría de los equipos combina prácticas
\cite{kuhrmann2022}.

\paragraph{Controles para la debilidad del ágil.}
Los problemas más frecuentes del ágil incluyen requisitos volátiles, clientes
que cambian de opinión, contratos de precio fijo y presión de plazos
\cite{meckenstock2024}. Para mitigarlos se establecen:

\begin{itemize}
  \item Requisitos de la LOPDP como historias obligatorias desde el sprint~1:
  aviso de privacidad, consentimiento explícito, uso de datos solo para el
  pedido, cifrado y derecho a eliminar la cuenta
  \cite{canedo2022,lopdp2021,reglamento2023}.
  \item Una Definición de Hecho que incluya prueba en Android de gama baja,
  funcionamiento con conexión intermitente y revisión del manejo de datos
  personales, para no descuidar los requisitos no funcionales
  \cite{schon2017}.
  \item Conciliación de transferencias como requisito obligatorio: ningún
  pedido pasa a cocina sin pago confirmado, o se muestra como
  ``pendiente de pago''.
  \item Priorización MoSCoW acordada por escrito con la dueña, porque el
  contrato de precio fijo es informal.
  \item Regla de cambios: las ideas nuevas (puntos, promociones, propuestas
  del hijo) entran al backlog, nunca al sprint en curso.
\end{itemize}

La Tabla~\ref{tab:moscow} presenta el alcance priorizado.

\begin{table}[htbp]
\centering
\caption{Alcance priorizado con MoSCoW.}
\label{tab:moscow}
\small
\begin{tabularx}{\textwidth}{@{}P{3.2cm}L@{}}
\toprule
\textbf{Prioridad} & \textbf{Funcionalidades} \\
\midrule
\emph{Must} (MVP) & Menú del día editable por la dueña, carrito y pedido, retiro o domicilio con dirección y referencia, pago por transferencia con carga de comprobante, panel de pedidos y conciliación, consentimiento de datos \\
\emph{Should} & Modo sin conexión (caché del menú y cola de pedidos), notificaciones de estado, historial de pedidos \\
\emph{Could} & Promociones por día o programa de puntos (se elige uno), compartir en redes \\
\emph{Won't} (esta versión) & Pago con tarjeta, múltiples locales \\
\bottomrule
\end{tabularx}
\end{table}

\subsection{Framework de IR por iteración}

En el modelo ágil, las cuatro actividades centrales de IR ---elicitación,
análisis y negociación, especificación y validación--- se repiten en cada
sprint a pequeña escala \cite{schon2017}, mientras que las transversales
---gestión de cambios, trazabilidad, priorización y comunicación con los
actores--- se ejecutan de forma continua durante las 14 semanas
\cite{kasauli2021}. La Tabla~\ref{tab:framework} detalla su ejecución.

% ---------- Tabla horizontal del framework ----------
\begin{landscape}
\begin{table}[p]
\centering
\caption{Actividades centrales y transversales del framework de IR en cada sprint.}
\label{tab:framework}
\scriptsize
\renewcommand{\arraystretch}{1.25}
\begin{tabularx}{\linewidth}{@{}P{1.6cm}P{2.2cm}LLLLL@{}}
\toprule
\textbf{Sprint (semanas)} & \textbf{Objetivo del sprint} & \textbf{Elicitación} & \textbf{Análisis y negociación} & \textbf{Especificación} & \textbf{Validación} & \textbf{Transversales destacadas} \\
\midrule
Sprint 0 (1--2) & Visión y backlog inicial &
Entrevista con Doña Rosa y su hijo; observación de una hora pico; revisión del cuaderno y los chats; encuesta corta a clientes &
Identificar actores, separar necesidades de deseos, primera negociación del MVP &
Visión del producto, mapa de historias, backlog inicial con criterios de aceptación, requisitos no funcionales (LOPDP, peso de la app, modo sin conexión) &
La dueña aprueba el MVP y la priorización MoSCoW por escrito &
Backlog como línea base; matriz de trazabilidad requisito $\rightarrow$ historia $\rightarrow$ prueba; acuerdo de comunicación (revisión quincenal, canal único de ideas) \\
\addlinespace
Sprint 1 (3--4) & Menú del día y catálogo &
Taller con la dueña sobre cómo cambia el menú y la disponibilidad diaria &
Simplificar la edición del menú para que la dueña la haga desde el teléfono &
Historias del menú y del panel de administración; prototipo en papel o Figma &
Demo en un Android de gama baja; la dueña publica un menú real &
Cambios de la demo registrados en el backlog; historias de privacidad (aviso y consentimiento) incluidas \\
\addlinespace
Sprint 2 (5--6) & Pedido, retiro o domicilio &
Prueba con 3--5 clientes frecuentes; entrevista al repartidor sobre direcciones &
Resolver el problema de direcciones: referencia, ubicación GPS opcional, teléfono &
Historias de carrito, pedido y entrega; reglas de validación de dirección &
Prueba de aceptación con pedidos simulados en hora pico &
Trazabilidad de los errores del cuaderno a las historias que los resuelven; ideas del hijo registradas y pospuestas \\
\addlinespace
Sprint 3 (7--8) & Pago por transferencia y conciliación (\textbf{MVP}) &
Sesión con la dueña sobre cómo verifica hoy las transferencias &
Negociar el flujo: pedido ``pendiente de pago'' hasta confirmar el comprobante &
Historias de carga de comprobante, estados del pedido y conciliación; requisitos de seguridad de datos &
Piloto real del MVP con clientes conocidos durante una semana &
Gestión de riesgos: pedidos sin pagar; revisión de cumplimiento LOPDP del MVP \\
\addlinespace
Sprint 4 (9--10) & Tolerancia a cortes y rendimiento &
Métricas del piloto: fallos de conexión, tiempos de carga, quejas &
Priorizar el modo sin conexión frente a nuevas funciones solicitadas &
Requisitos no funcionales medibles (p.\,ej., app menor a 15\,MB, menú visible sin conexión) &
Pruebas en teléfonos de gama baja con datos móviles intermitentes &
Repriorización del backlog con datos del piloto; actualización de la trazabilidad \\
\addlinespace
Sprint 5 (11--12) & Fidelización y notificaciones &
Taller con la dueña y el hijo: puntos o promociones por día &
Negociación: se elige una sola opción según presupuesto y plazo restante &
Historias de la opción elegida y notificaciones de estado del pedido &
Revisión del sprint con la familia y algunos clientes &
Control de cambios: la opción descartada queda en el backlog para después de las fiestas \\
\addlinespace
Sprint 6 (13--14) & Estabilización para las fiestas &
Revisión de defectos reportados; previsión de demanda (el doble de pedidos) &
Decidir qué se congela: no entran funciones nuevas &
Ajustes finales; manual breve para la dueña y aviso de privacidad definitivo &
Prueba de carga con el doble de pedidos; aceptación final de Doña Rosa &
Línea base final de requisitos; backlog posterior a las fiestas (pago con tarjeta, opción de fidelización descartada) \\
\bottomrule
\end{tabularx}
\end{table}
\end{landscape}

% ---------------------------------------------------------------------
\section{Discusión}

Los resultados coinciden con la literatura. Otavo Piraban y Pardo Moreno
\cite{otavo2024} emplean también un cuadro comparativo con criterios del
proyecto para elegir el proceso de una aplicación de pagos electrónicos, y
optan por un enfoque iterativo ante requisitos cambiantes y normas que
evolucionan. Mex Alvarez \emph{et al.} \cite{mex2025} observaron que la
cascada solo destaca en gestión de requisitos y muestra rigidez ante el
cambio, lo que respalda su baja valoración en C1, C3 y C6.

La elección del ágil no ignora sus riesgos. Meckenstock \cite{meckenstock2024}
identifica 90 problemas del desarrollo ágil, entre los que destacan la
presión de entrega y el comportamiento cambiante del cliente; por eso se
añadieron la regla de cambios y la priorización MoSCoW por escrito. En
cuanto a la privacidad, Canedo \emph{et al.} \cite{canedo2022} muestran que,
tras la entrada en vigor de una ley de datos similar a la LOPDP, la mayoría
de los equipos ágiles especifica la privacidad con historias de usuario y
considera que los requisitos son el área más afectada, lo que justifica
tratarla como historias obligatorias desde el sprint~1. La
Tabla~\ref{tab:respaldo} resume la relación entre cada decisión y la
evidencia que la respalda.

\begin{table}[htbp]
\centering
\caption{Respaldo de las decisiones en la literatura científica y la normativa.}
\label{tab:respaldo}
\footnotesize
\begin{tabularx}{\textwidth}{@{}P{4.3cm}LC{1.3cm}@{}}
\toprule
\textbf{Decisión de la actividad} & \textbf{Qué aporta la evidencia} & \textbf{Ref.} \\
\midrule
Matriz comparativa con criterios ponderados & Cuadro comparativo de modelos según los criterios del proyecto para una app de pagos electrónicos; se elige un enfoque iterativo por requisitos cambiantes & \cite{otavo2024} \\
Cascada con baja valoración en volatilidad, plazo y gestión del cambio & En proyectos de 3 meses, Scrum superó a cascada en gestión del cambio, planificación y adaptabilidad & \cite{mex2025} \\
IR repetida en cada sprint y participación diaria de la dueña & Las actividades de IR se repiten en cada iteración; la participación continua del cliente es factor crítico de éxito & \cite{schon2017} \\
Requisitos no funcionales en la Definición de Hecho & El descuido de los requisitos no funcionales es un desafío típico del ágil que provoca retrabajo & \cite{schon2017} \\
Privacidad como historias obligatorias desde el sprint 1 & Los equipos ágiles especifican la privacidad con historias de usuario tras la ley de datos & \cite{canedo2022} \\
Controles contra las debilidades del ágil & Requisitos volátiles, cliente cambiante, precio fijo y presión de plazos son problemas frecuentes & \cite{meckenstock2024} \\
Scrum ligero adaptado (no puro) & Menos del 15\,\% trabaja de forma puramente ágil o tradicional & \cite{kuhrmann2022} \\
Gestión de cambios y trazabilidad como transversales & 24 desafíos de IR en contextos ágiles, entre ellos gestionar el conocimiento de requisitos que cambian & \cite{kasauli2021} \\
Criterio C5 y requisitos de consentimiento, finalidad y eliminación & Marco legal obligatorio para tratar datos de clientes en Ecuador & \cite{lopdp2021,reglamento2023} \\
\bottomrule
\end{tabularx}
\end{table}

Como limitación, la valoración de la matriz es cualitativa y depende del
juicio del analista; otro equipo podría asignar puntajes distintos, aunque es
poco probable que cambie el orden de los modelos dada la diferencia de 15
puntos entre ágil y cascada. Además, la LOPDP fue reformada en su
artículo~43 (notificación de brechas) por la Ley Orgánica para el
Fortalecimiento de la Ciberseguridad de mayo de 2026, por lo que conviene
verificar la normativa vigente al implementar el sistema.

% ---------------------------------------------------------------------
\section{Conclusiones}

\begin{itemize}
  \item La alta volatilidad de los requisitos y la gran disponibilidad de la
  dueña hacen del ágil el modelo más adecuado (31/35); la cascada es el menos
  adecuado (16/35) porque congela los requisitos y entrega todo al final.
  \item La fecha fija se protege con alcance variable: el MVP (menú, pedido,
  entrega, pago por transferencia y consentimiento de datos) está listo en la
  semana~8, y las seis semanas restantes mejoran la calidad y añaden
  funciones según su prioridad.
  \item La debilidad del ágil en regulación se compensa incluyendo la LOPDP y
  la conciliación de pagos como requisitos obligatorios en la Definición de
  Hecho desde el primer sprint.
  \item La gestión de cambios es la actividad transversal crítica: toda idea
  nueva entra al backlog, nunca al sprint en curso, para que la volatilidad no
  ponga en riesgo la fecha de las fiestas.
\end{itemize}

% ---------------------------------------------------------------------
% Referencias en formato IEEE desde el archivo referencias.bib
\renewcommand{\refname}{Referencias}
{\small
\bibliographystyle{IEEEtran}
\bibliography{referencias}}

\end{document}
